\paragraph{二尺度反混叠与离散曲率刚性的统一误差预算}
设有限缺陷指数谱像写为
\[
u_n(\Delta)=4\pi\sum_{j=1}^{\kappa}\frac{m_j\delta_j}{1+\delta_j}
\exp\!\bigl(-(1+\delta_j)n\Delta\bigr)\exp\!\bigl(-i\gamma_j n\Delta\bigr),
\qquad n\ge 0,
\]
其中 \(\kappa\ge 1\)、\(\delta_j>0\)、\(\gamma_j\in\mathbb R\)、\(m_j>0\)。
记单模态谱像
\[
z_j(\Delta):=\exp\!\bigl(-(1+\delta_j)\Delta\bigr)\exp(-i\gamma_j\Delta)\in\mathbb D.
\]

\begin{theorem}[二尺度共动指数谱像的无混叠可识别性与 Diophantine 相位门槛]\label{thm:xi-two-scale-comoving-spectrum-diophantine-threshold}
\leanverified{paper\_xi\_two\_scale\_comoving\_spectrum\_diophantine\_threshold}
取两步长 \(\Delta_1,\Delta_2>0\)，并设
\[
\alpha:=\frac{\Delta_1}{\Delta_2}\notin\mathbb Q.
\]
则成立：
\begin{enumerate}
  \item 映射
  \[
  (\gamma,\delta)\longmapsto \bigl(z(\Delta_1),z(\Delta_2)\bigr)
  \]
  在 \(\mathbb R\times(0,\infty)\) 上单射，故 \(\bigl(z(\Delta_1),z(\Delta_2)\bigr)\) 唯一决定 \((\gamma,\delta)\)；
  \item 进一步设 \(\alpha\) 满足 \((c,\nu)\)-Diophantine 条件：存在 \(c>0\)、\(\nu\ge2\)，对任意既约有理数 \(p/q\)（\(q\ge1\)）有
  \[
  \left|\alpha-\frac{p}{q}\right|\ge\frac{c}{q^\nu}.
  \]
  不失一般性取 \(\Delta_1\le\Delta_2\)。若 \(|\gamma|\le T\)，并定义观测相位误差
  \[
  \eta:=\max_{r=1,2}\operatorname{dist}\!\bigl(\arg z^{\mathrm{obs}}(\Delta_r)-\arg z(\Delta_r),\,2\pi\mathbb Z\bigr),
  \]
  当
  \[
  \eta<
  \frac{\pi c}{
  \bigl(\lceil T\Delta_2/(2\pi)\rceil+1\bigr)^{\nu-1}
  }
  \]
  时，\(\gamma\) 的整数提升在区间 \([-T,T]\) 内唯一，且任意一致重构 \(\gamma^{\mathrm{rec}}\) 满足
  \[
  |\gamma^{\mathrm{rec}}-\gamma|
  \le
  \frac{\eta}{\Delta_1}+\frac{\eta}{\Delta_2}.
  \]
  同时
  \[
  |\delta^{\mathrm{rec}}-\delta|
  \le
  \min_{r=1,2}
  \frac{1}{\Delta_r}
  \left|
  \log|z^{\mathrm{obs}}(\Delta_r)|-\log|z(\Delta_r)|
  \right|.
  \]
\end{enumerate}
\end{theorem}
\begin{proof}
由
\[
|z(\Delta_r)|=\exp\!\bigl(-(1+\delta)\Delta_r\bigr)
\]
可得
\[
\delta=-\frac{1}{\Delta_r}\log|z(\Delta_r)|-1,
\]
故 \(\delta\) 由任一通道模长唯一确定。
若两组参数 \((\gamma,\delta)\)、\((\gamma',\delta)\) 诱导同一双通道相位，则
\[
e^{-i(\gamma-\gamma')\Delta_r}=1,\qquad r=1,2,
\]
即
\[
(\gamma-\gamma')\Delta_r\in2\pi\mathbb Z.
\]
从而存在整数 \(k_1,k_2\) 使
\[
\gamma-\gamma'=\frac{2\pi k_1}{\Delta_1}=\frac{2\pi k_2}{\Delta_2},
\]
故 \(\alpha=\Delta_1/\Delta_2=k_1/k_2\)。由 \(\alpha\notin\mathbb Q\) 得 \(k_1=k_2=0\)，即 \(\gamma=\gamma'\)。单射性成立。

对误差门槛部分，取相位代表元 \(\theta_r\in\mathbb R\) 满足
\[
e^{-i\theta_r}=\frac{z(\Delta_r)}{|z(\Delta_r)|},
\qquad
\theta_r=-\gamma\Delta_r.
\]
观测相位写为
\[
\arg z^{\mathrm{obs}}(\Delta_r)=\theta_r+e_r\pmod{2\pi},
\qquad |e_r|\le\eta.
\]
任意重构提升整数 \(k_r\in\mathbb Z\) 给出
\[
\gamma^{\mathrm{rec}}\Delta_r=-\theta_r-e_r+2\pi k_r.
\]
于是
\[
\Delta_r(\gamma^{\mathrm{rec}}-\gamma)=2\pi k_r-e_r.
\]
若发生错误提升，则 \((k_1,k_2)\neq(0,0)\)。
将上式分别除以 \(\Delta_1,\Delta_2\) 并相减：
\[
\left|
\frac{k_1}{\Delta_1}-\frac{k_2}{\Delta_2}
\right|
\le
\frac{\eta}{2\pi}\left(\frac1{\Delta_1}+\frac1{\Delta_2}\right).
\]
若 \(k_2=0\)，则由上式与 \(\Delta_1\le\Delta_2\) 得
\[
\frac{|k_1|}{\Delta_1}
\le
\frac{\eta}{2\pi}\left(\frac1{\Delta_1}+\frac1{\Delta_2}\right)
\le
\frac{\eta}{\pi\Delta_1}
<
\frac1{\Delta_1},
\]
故 \(k_1=0\)，与错误提升矛盾，因此必有 \(k_2\neq0\)。
于是
\[
\left|\frac{k_1}{k_2}-\alpha\right|
\le
\frac{\eta}{2\pi|k_2|}(1+\alpha).
\]
由 \(\Delta_1\le\Delta_2\) 得 \(\alpha\le1\)，故右侧不超过 \(\eta/(\pi|k_2|)\)。
另一方面，若候选频率 \(\gamma^{\mathrm{rec}}\in[-T,T]\)，则
\[
|k_2|\le \left\lceil\frac{T\Delta_2}{2\pi}\right\rceil+1=:Q_T.
\]
Diophantine 下界给出
\[
\left|\frac{k_1}{k_2}-\alpha\right|
\ge
\frac{c}{|k_2|^\nu}
\ge
\frac{c}{Q_T^\nu}.
\]
若
\[
\eta<\frac{\pi c}{Q_T^{\nu-1}},
\]
则与上界矛盾，故不存在非零 \((k_1,k_2)\)，提升唯一。

在唯一提升情形 \(k_r=0\) 下，
\[
|\gamma^{\mathrm{rec}}-\gamma|
=\frac{|e_r|}{\Delta_r}\le\frac{\eta}{\Delta_r},\qquad r=1,2,
\]
从而
\[
|\gamma^{\mathrm{rec}}-\gamma|
\le
\frac{\eta}{\Delta_1}+\frac{\eta}{\Delta_2}.
\]
\(\delta\) 的误差界由
\[
\delta=-\frac1{\Delta_r}\log|z(\Delta_r)|-1
\]
逐通道直接传播并取最小值得证。证毕。
\end{proof}

\begin{corollary}[黄金比例步长比值的统一门槛常数]\label{cor:xi-golden-ratio-stepsize-unaliasing-threshold}
若
\[
\alpha=\varphi:=\frac{1+\sqrt5}{2},
\]
则由 Hurwitz--Markov 极值性质，对任意既约有理数 \(p/q\) 有
\[
\left|\varphi-\frac{p}{q}\right|\ge\frac{1}{\sqrt5\,q^2}.
\]
因此定理 \ref{thm:xi-two-scale-comoving-spectrum-diophantine-threshold} 中可取
\[
(c,\nu)=\left(\frac1{\sqrt5},\,2\right),
\]
并得到统一相位容忍阈值
\[
\eta<
\frac{\pi}{
\sqrt5\left(\left\lceil T\Delta_2/(2\pi)\right\rceil+1\right)
}
\asymp T^{-1}.
\]
\end{corollary}
\begin{proof}
将
\(
c=1/\sqrt5,\ \nu=2
\)
代入定理 \ref{thm:xi-two-scale-comoving-spectrum-diophantine-threshold} 的门槛式即可。证毕。
\end{proof}

\begin{theorem}[离散曲率控制下的近单模态刚性]\label{thm:xi-discrete-curvature-near-unimodal-rigidity}
设有限谱
\[
\{a_j\}_{j=1}^{\kappa}\subset[1,\infty),\qquad w_j>0,
\]
并定义
\[
Z(u):=\sum_{j=1}^{\kappa}w_j e^{-ua_j},\qquad u\ge0.
\]
令倾斜分布
\[
\mathbb P_u(A=a_j):=\frac{w_j e^{-ua_j}}{Z(u)}.
\]
则
\[
\partial_u\log Z(u)=-\mathbb E_u[A],\qquad
\partial_u^2\log Z(u)=\Var_u(A)\ge0.
\]
给定 \(\Delta>0\)，定义离散曲率
\[
\kappa_Z(u;\Delta):=\log Z(u+\Delta)-2\log Z(u)+\log Z(u-\Delta),
\qquad u\ge\Delta.
\]
则成立：
\begin{enumerate}
  \item 精确积分表示
  \[
  \kappa_Z(u;\Delta)=
  \int_{-\Delta}^{\Delta}(\Delta-|s|)\Var_{u+s}(A)\,\dd s;
  \]
  \item 若 \(\kappa_Z(u;\Delta)\le\varepsilon\)，则存在 \(u^\ast\in[u-\Delta,u+\Delta]\) 使
  \[
  \Var_{u^\ast}(A)\le\frac{2\varepsilon}{\Delta^2}.
  \]
  因而对任意 \(r>0\) 有
  \[
  \mathbb P_{u^\ast}\!\bigl(|A-\mathbb E_{u^\ast}A|\ge r\bigr)
  \le
  \frac{2\varepsilon}{\Delta^2r^2}.
  \]
\end{enumerate}
\end{theorem}
\begin{proof}
对任意 \(f\in C^2\) 有恒等式
\[
f(u+\Delta)-2f(u)+f(u-\Delta)
=
\int_{-\Delta}^{\Delta}(\Delta-|s|)f''(u+s)\,\dd s.
\]
取 \(f=\log Z\) 即得第一条。

若对一切 \(s\in[-\Delta,\Delta]\) 都有
\(
\Var_{u+s}(A)>2\varepsilon/\Delta^2
\)，
则由第一条与
\[
\int_{-\Delta}^{\Delta}(\Delta-|s|)\,\dd s=\Delta^2
\]
推出 \(\kappa_Z(u;\Delta)>2\varepsilon\)，与假设矛盾，故存在 \(u^\ast\) 满足所述方差上界。
尾界由 Chebyshev 不等式直接给出。证毕。
\end{proof}

\begin{theorem}[二尺度反混叠与离散曲率刚性的全链路停止证书]\label{thm:xi-two-scale-unaliasing-discrete-curvature-stop-certificate}
在定理 \ref{thm:xi-two-scale-comoving-spectrum-diophantine-threshold} 的模型下，设两通道步长比值 \(\alpha=\Delta_1/\Delta_2\) 满足 \((c,\nu)\)-Diophantine 条件，并设同一缺陷谱诱导配分函数 \(Z\)。
若在某尺度参数 \(u\) 上有
\[
\kappa_Z(u;\Delta_r)\le\varepsilon_r,\qquad r=1,2,
\]
则存在 \(u_r^\ast\in[u-\Delta_r,u+\Delta_r]\) 使
\[
\Var_{u_r^\ast}(A)\le\frac{2\varepsilon_r}{\Delta_r^2}.
\]
进一步假设重构流程给出残差 \(\mathcal R_r\) 与相位误差 \(\eta_r\)，并满足可审计放大模型
\[
\eta_r\le \Lambda_r\!\bigl(\Var_{u_r^\ast}(A)\bigr)\,\mathcal R_r,
\]
其中 \(\Lambda_r\) 对其自变量单调不减。
则有
\[
\eta_r\le
\Lambda_r\!\left(\frac{2\varepsilon_r}{\Delta_r^2}\right)\mathcal R_r.
\]
记 \(\eta:=\max\{\eta_1,\eta_2\}\)。若
\[
\eta<
\frac{\pi c}{
\bigl(\lceil T\Delta_2/(2\pi)\rceil+1\bigr)^{\nu-1}
},
\]
则所有 \(|\gamma_j|\le T\) 模态的整数提升唯一，且
\[
|\gamma_j^{\mathrm{rec}}-\gamma_j|
\le
\frac{\eta_1}{\Delta_1}+\frac{\eta_2}{\Delta_2},
\]
\[
|\delta_j^{\mathrm{rec}}-\delta_j|
\le
\min_{r=1,2}
\frac{1}{\Delta_r}
\left|
\log|z_j^{\mathrm{obs}}(\Delta_r)|-\log|z_j(\Delta_r)|
\right|.
\]
因此，一旦曲率阈值与 Diophantine 相位阈值同时满足，即得到可审计停止条件：后续递归仅改变常数项，不再改变无混叠可识别性与误差阶。
\end{theorem}
\begin{proof}
前半部分由定理 \ref{thm:xi-discrete-curvature-near-unimodal-rigidity} 逐通道应用得到。
单调性假设给出
\[
\eta_r\le
\Lambda_r\!\left(\frac{2\varepsilon_r}{\Delta_r^2}\right)\mathcal R_r.
\]
再将 \(\eta=\max\{\eta_1,\eta_2\}\) 代入定理 \ref{thm:xi-two-scale-comoving-spectrum-diophantine-threshold} 即得提升唯一性与 \((\gamma_j,\delta_j)\) 的显式误差界。
由于该界的主阶由阈值结构决定，阈值满足后继续递归不改变识别性与误差阶，仅可能改善前常数，故停止判据成立。证毕。
\end{proof}

\endinput
